\subsection{Laplacianizer moduli}

\begin{theorem}[Moduli of laplacianizers (orbifold)]\label{thm:laplacianizer-moduli-orbifold}\notready\uses{thm:settling-hodge-iso}
\emph{Open frontier, stated but not formalized.} The continuous moduli of
laplacianizers: the space of compatible inner products on the cochain modules,
an orbifold quotient by the gauge group, of which the discrete algebraic
Laplacianizer (\cref{thm:settling-hodge-iso}) is the inner-product-free
projection. Stating it needs the real layer: inner products and the orbifold
structure. Red until a constructive real layer carries the continuous moduli.
\end{theorem}

\begin{theorem}[Cellular sheaf Laplacian as a moduli case (Hansen--Ghrist subsumption)]\label{thm:cellular-sheaf-laplacian-moduli-case}\notready\uses{thm:laplacianizer-moduli-orbifold}
\emph{Open frontier, stated but not formalized.} The real-valued moduli of
laplacianizations (\cref{thm:laplacianizer-moduli-orbifold}) is a real-analytic
orbifold stratified by spectral multipartition with category-independent
codimension $\sum_{|B|\ge2}(|B|-1)$; the cellular sheaf Laplacian of
Hansen--Ghrist and Ayzenberg--Magai is the cellular finite-dimensional case,
its realizability obstruction the univariate IEP-G program (first non-trivial
obstruction $P_3$, Barrett--Fallat). The cellular sheaf-Laplacian program is
therefore a stratum-theoretic special case of the moduli claim. Stating it
needs the real layer: Hilbert cochains, SPD-cone moduli, and eigenvalue strata.
Red until a constructive real layer carries the inner-product moduli and its
stratification.
\end{theorem}
